\paragraph{Alcance de la consulta implementada.}\label{sec:diseno-bitacora}
El incremento de consulta desarrolla la selección cronológica prevista en la
\cref{tab:cu-17}: un Owner autenticado puede solicitar un intervalo temporal y
filtros exactos de actor registrado, operación y recurso. La consulta exige
permisos vigentes incluso cuando no existen coincidencias. Las páginas conservan
el máximo de secuencia de su primera lectura, por lo que los anexados posteriores
no se incorporan a esa misma navegación aunque tengan una fecha anterior.

El intervalo incluye su límite inicial y excluye el final, tiene como máximo
366 días y usa instantes UTC con precisión de nanosegundos. El orden se define
por segundos, nanosegundos y secuencia; esta última resuelve eventos con la misma
marca temporal. Cada página contiene hasta cien eventos completos. El cursor
vincula intervalo, filtros, última posición y máximo de secuencia: selecciona
una continuación, pero no acredita identidad, permisos ni integridad de la cadena.

La lectura de actividad y la verificación criptográfica tienen responsabilidades
distintas. La primera selecciona registros autorizados y registra su propia
consulta; la segunda recalcula los eslabones. Una página no es una verificación
del historial completo, y ninguna de ambas demuestra por sí sola que no se haya
retirado el extremo final de una cadena válida.

Las \cref{tab:cu-17,tab:rf-20} mantienen su alcance aprobado. El actor histórico
es el texto que ya forma parte del evento, habitualmente un correo; no se presenta
como identificador estable de la cuenta. Los eventos actuales no incorporan la
dirección IP exigida por esas tablas. Migrar la consulta no autoriza reconstruir
identidades o direcciones ausentes ni modificar el contenido encadenado. El
criterio de registro en menos de 500 milisegundos y la cobertura de todas las
operaciones requieren evidencia propia. La consulta parcial no declara cumplido
el caso de uso completo.
